\documentclass[11pt]{article}
\usepackage[margin=2.6cm]{geometry}
\usepackage{amsmath,amssymb,amsthm}
\usepackage{booktabs}

\newtheorem{theorem}{Theorem}
\newtheorem{proposition}{Proposition}
\newtheorem{corollary}{Corollary}
\newtheorem{remark}{Remark}

\newcommand{\Q}{\mathbf{Q}}

\title{A tenth-discriminant bound from Guo's zero gap:\\
an explicit extraction via Kimura's Proposition~9}
\author{George A. D. Morris \and Ozzie (AI research assistant, Muse, Meta)}
\date{30 October 2026}

\begin{document}
\maketitle

\begin{abstract}
In response to Terence Tao's ``tenth discriminant'' challenge (20 October
2026), which asks for explicit, optimised upper bounds on a hypothetical
tenth class-number-one imaginary quadratic discriminant derived from the
recent proofs excluding Landau--Siegel zeros, we extract such a bound from
the zero-gap formulation of Kimura, using the explicit zero gap of Guo in
place of Kimura's numerically certified grid constants. Guo's Theorem~1.1
gives $(1-\beta)\log q > 1/100$ for every primitive quadratic Dirichlet
character of conductor $q \ge 3$ and every real zero $\beta$; feeding the
zero gap $(c,T) = (1/100, T)$ into Kimura's Proposition~9 and the class
number formula, and certifying the resulting elementary inequality in
interval arithmetic, we obtain
\[
h(-d) = 1 \ \Longrightarrow\ d \le 32{,}605{,}775,
\]
and, with the stronger constant $c = 1/64$ of Guo's Remark~1.2,
$d \le 12{,}609{,}991$. These improve every previously posted extraction
through Kimura's Proposition~9 (Rudolph: $d < 189{,}348{,}540$; Murphy's
chains: $6.47\times 10^7$ and $1.23\times 10^8$) by a factor of roughly
$2$--$15$. We are explicit that the bound is not competitive with the
determinant-method bounds ($d \le 331$) or the quasi-Riemann-hypothesis
bounds ($d \le 653$) posted in the same challenge; its interest is as the
best bound obtainable from the zero-gap route, by an argument whose only
non-elementary inputs are two cited, fully explicit theorems.
\end{abstract}

\section{Introduction}

The nine Heegner discriminants $d = 1, 2, 3, 7, 11, 19, 43, 67, 163$ are
exactly the positive integers for which $\Q(\sqrt{-d})$ has class number
one (Heegner--Stark--Baker). If $d$ is a Heegner number and $\chi$ is the
associated quadratic character, the class number formula gives
$L(1,\chi) = \pi/\sqrt{d}$, and a short argument produces a Landau--Siegel
zero $\beta$ of $L(s,\chi)$ with $1-\beta \asymp 1/\sqrt{d}$. Consequently,
any sufficiently explicit theorem excluding real zeros of quadratic
Dirichlet $L$-functions close to $s=1$ yields an explicit upper bound on a
hypothetical tenth discriminant. Making this effective by Goldfeld's
method gave $d \le e^{7000}$, improved by Oesterl\'e to $d \le e^{55}$.

In October 2026, OpenAI's manuscript \cite{openai} asserting the uniform
exclusion of Landau--Siegel zeros (and, in its family~003, a zero-free
region $\Re s > 7/8$) was followed within two weeks by independent proofs
due to Banks \cite{banks}, Kerr \cite{kerr}, Green \cite{green} and
Kimura \cite{kimura}, among others. On 20 October 2026 Tao launched an
open challenge \cite{tao}: for each such proof, or by a novel argument,
\emph{use an AI tool to work out an effective upper bound for a
hypothetical tenth discriminant, without using Heegner's methods, and make
best efforts to optimise the constant}, with contributions posted in the
comments of the challenge post, AI assistance disclosed, and attribution
to prior contributions.

The strongest bounds posted in the challenge to date come from the
determinant method (Chung/rStation agents: $D \le 331$, 28 October 2026;
Murphy: $q < 565$; twinheterodox: $q < 570$) and from evaluating the
Chowla--Selberg formula at the fixed point $s = 7/8$ of the quasi-Riemann
hypothesis (LemmaLab: $d \le 653$). A separate, simpler family of
contributions extracts bounds from \emph{zero gaps} via Kimura's
Proposition~9 \cite{kimura}: Rudolph \cite{tao} obtained
$|D| < e^{19.0591} < 189{,}348{,}540$ using Kimura's grid constant
$c = 0.004401$, and Murphy's survey records two further chains of this
type giving $d < 6.47\times 10^7$ and $d < 1.23\times 10^8$.

The present note improves the zero-gap family. Kimura's grid constants
grow from $5.811\times 10^{-4}$ at small conductors to
$5.965\times 10^{-3}$ for $\log q \ge 35.0131$. Guo \cite{guo} has since
proved, by an elementary argument, the uniform bound
$(1-\beta)\log q > 1/100$ for \emph{all} conductors $q \ge 3$ (and
$1/64$ in his Remark~1.2), which is stronger than every constant in
Kimura's table in the range relevant here. Substituting Guo's constant
into Kimura's conversion and optimising the free parameters yields the
bounds in the abstract. The derivation uses no finite conductor-range
computation: unlike Banks's all-conductor corollary, which rests on the
finite-range computation of Lu--Zaman--Zhao, Guo's theorem covers all
$q \ge 3$ directly, and Kimura's Proposition~9 is proved in Lean in
\cite{kimura}. The price of this cleanliness is tightness: the zero-gap
to class-number conversion is lossy, and our bound does not approach the
determinant-method leaders. Tao's challenge explicitly welcomes
per-approach bounds of this kind.

\section{The three inputs}

Throughout, $d \ge 3$ denotes the absolute value of a fundamental
discriminant $D_K < 0$ of an imaginary quadratic field $K$, $\chi$ is the
primitive quadratic character of conductor $d$, and $h(-d)$ is the class
number.

\begin{proposition}[Guo \cite{guo}, Theorem 1.1 and Remark 1.2]
\label{prop:guo}
For every primitive quadratic Dirichlet character $\chi$ of conductor
$q \ge 3$, every real zero $\beta \in (0,1)$ of $L(s,\chi)$ satisfies
$(1-\beta)\log q > 1/100$. The same argument yields the constant $1/64$ in
place of $1/100$ \textup{(}Remark 1.2, verified in Appendix A of
\cite{guo}\textup{)}.
\end{proposition}

Kimura \cite{kimura} calls $(c,T)$ a \emph{zero gap} if
$(1-\beta)\ell \ge c$ for every primitive real character of conductor
$q \ge 3$ with $\ell = \log q \ge T$ and every real zero $\beta$.
Proposition~\ref{prop:guo} says precisely that the zero gap $(c,T)$ holds
with $c = 1/100$ (resp.\ $c = 1/64$) for \emph{every} $T$.

\begin{proposition}[Kimura \cite{kimura}, Proposition 9]
\label{prop:kimura}
Let $c > 0$, $T \ge 2$, $c/T \le 1/2$, and suppose the zero gap $(c,T)$
holds. Let $A > 0$ with $AT \ge e$, let $V \ge 1$ be an integer, and put
$\sigma_T = 1 - c/T$, $P = c\bigl(1 + \log(2VAT)/T\bigr)$, and
\[
\varepsilon
= e^{c}\Bigl(\frac{(AT)^{c/T} V^{c/T}}{cA}
+ \frac{4}{(AT)^{\sigma_T}}\Bigr)
+ \frac{e^{P}}{cA}
+ \frac{4e^{P}}{V^{\sigma_T}}
+ \frac{e^{P}}{AT}.
\]
If $\varepsilon < 1$, then every primitive real $\chi$ of conductor
$q \ge 3$ with $\log q \ge T$ satisfies
\[
\sqrt{q}\, L(1,\chi) \;\ge\; (1-\varepsilon)\, c\, e^{-P}\, e^{T/2}/T.
\]
\end{proposition}

\begin{proposition}[Class number formula; \cite{kimura}, eq.\ (4)]
\label{prop:cnf}
$2\pi h(-d) = w_K \sqrt{d}\, L(1,\chi)$, where $w_K \ge 2$ is the number of
roots of unity in $K$. Hence $\pi h(-d) \ge \sqrt{d}\, L(1,\chi)$.
\end{proposition}

\section{Main result}

\begin{theorem}
\label{thm:main}
If $h(-d) = 1$ then $d \le 32{,}605{,}775$.
\end{theorem}

\begin{corollary}
\label{cor:variant}
With $c = 1/64$ \textup{(}Proposition~\ref{prop:guo}, second
statement\textup{)}, if $h(-d) = 1$ then $d \le 12{,}609{,}991$.
\end{corollary}

\begin{proof}[Proof of Theorem~\ref{thm:main}]
Suppose $h(-d) = 1$ and $\log d \ge T := 17.3$. By
Proposition~\ref{prop:guo} with $c = 1/100$, the zero gap $(c, T)$ holds.
Apply Proposition~\ref{prop:kimura} with $A = 10^5$ and $V = 10^4$ (so
$AT = 1.73\times 10^6 \ge e$ and $c/T \le 1/2$). Interval arithmetic
(Section~\ref{sec:verify}) certifies
$\varepsilon \le 0.00246309 < 1$ and
\[
\sqrt{d}\, L(1,\chi) \;\ge\; (1-\varepsilon)\, c\, e^{-P}\, e^{T/2}/T
\;\ge\; 3.21436540 \;>\; \pi.
\]
But Proposition~\ref{prop:cnf} with $h(-d) = 1$ gives
$\sqrt{d}\, L(1,\chi) \le \pi h(-d) = \pi$: a contradiction. Hence
$\log d < 17.3$, i.e.\ $d < e^{17.3} = 32{,}605{,}775.72\ldots$, so
$d \le 32{,}605{,}775$.
\end{proof}

\begin{proof}[Proof of Corollary~\ref{cor:variant}]
Identical, with $c = 1/64$, $T = 16.35$, $A = 10^5$, $V = 10^4$:
interval arithmetic certifies $\varepsilon \le 0.00175279$ and
$(1-\varepsilon) c e^{-P} e^{T/2}/T \ge 3.25884426 > \pi$, whence
$d < e^{16.35} = 12{,}609{,}991.07\ldots$
\end{proof}

\section{Computation and verification}
\label{sec:verify}

The only computation in the proof is the evaluation of the elementary
expressions in Proposition~\ref{prop:kimura}. It is performed by a short
Python script (\texttt{compute\_bound.py}, distributed with this note),
in three parts.

\paragraph{Validation of the machinery.}
Before optimising, we validated our implementation of $\varepsilon$ and
$P$ against Kimura's Theorem~10 \cite{kimura}, which uses
$A = 10^6$, $V = 10^4$ and his grid constants. Our computed log-margins
$\log\bigl((1-\varepsilon)c_N e^{-P} e^{T_N/2}/(T_N \pi)\bigr)$ are:

\begin{center}
\begin{tabular}{lccccc}
\toprule
$N$ & 100 & 200 & 500 & 1000 & 2000 \\
\midrule
margin & 0.01990 & 0.03129 & 0.00925 & 0.02373 & 0.07199 \\
\bottomrule
\end{tabular}
\end{center}

Kimura states that these margins lie between $0.009$ ($N = 500$) and
$0.072$ ($N = 2000$) and that $\varepsilon < 7.8\times 10^{-4}$; our
implementation returns $\varepsilon \le 7.73\times 10^{-4}$ across the
table, matching both. As a second, independent cross-check, running the
same machinery with Kimura's grid constant $c = 0.004401$ returns the
threshold $T = 19.060$, $e^T = 189{,}519{,}030$, agreeing with Rudolph's
independently extracted $e^{19.0591} \approx 1.89\times 10^8$ \cite{tao}.

\paragraph{Optimisation.}
A grid search over $A \in [10^3, 10^8]$, $V \in [10^2, 10^6]$,
$T \in [12, 22]$ finds float-arithmetic thresholds $T = 17.250$ for
$c = 1/100$ (at $A = 10^5$, $V = 10^4$) and $T = 16.270$ for $c = 1/64$
(at $A = 5.6\times 10^4$, $V = 10^4$). The certified values
$T = 17.3$ and $T = 16.35$ are round choices just past the optima,
trading negligible tightness for comfortable margins.

\paragraph{Certified values.}
The final inequalities were re-evaluated in rigorous interval
arithmetic (mpmath \texttt{iv}, 50-digit working precision). For the
headline parameters, $\sigma_T = 0.99942197$, $P = 0.02402724$, and the
five terms of $\varepsilon$ are
$1.0239\times 10^{-3}$, $2.3548\times 10^{-6}$, $1.0243\times 10^{-3}$,
$4.1191\times 10^{-4}$, $5.9209\times 10^{-7}$, summing to
$\varepsilon \le 0.00246309$; the certified lower bound is
$(1-\varepsilon) c e^{-P} e^{T/2}/T \ge 3.21436540 > \pi = 3.14159265\ldots$,
a margin of $+0.0728$. For the variant, $\sigma_T = 0.99904434$,
$P = 0.03876208$, $\varepsilon \le 0.00175279$, and the certified lower
bound is $3.25884426$, a margin of $+0.1173$. The margins are two orders
of magnitude larger than any conceivable rounding artefact, and the
certified digits are stable under precision changes.

\section{Placement and limitations}

\begin{center}
\begin{tabular}{lll}
\toprule
Bound & Route & Source \\
\midrule
$D \le 331$ & determinant & Chung/rStation, 28 Oct 2026 \\
$d \le 653$ & Chowla--Selberg at $7/8$ & LemmaLab \\
$q < 565$ & determinant & Murphy \\
$d < 1.89\times 10^8$ & Kimura Prop.~9, $c = 0.004401$ & Rudolph \\
$d < 6.47\times 10^7$ & Kimura zero-gap chain 58 & Murphy (survey) \\
$d < 1.23\times 10^8$ & Kimura zero-gap chain 59 & Murphy (survey) \\
$d \le 32{,}605{,}775$ & Kimura Prop.~9, Guo $c = 1/100$ & this note \\
$d \le 12{,}609{,}991$ & Kimura Prop.~9, Guo $c = 1/64$ & this note \\
\bottomrule
\end{tabular}
\end{center}

Three limitations deserve emphasis. First, this is a \emph{bound}, not a
classification: we do not enumerate the range $d < e^{17.3}$ (Kimura's
Theorem~2 machinery, or the sieves of Larsson, Rudolph and Murphy, address
that separate step). Second, the argument depends on Guo's preprint
\cite{guo}, which at the time of writing is an unrefereed v1 and, unlike
Kimura's results, is not formally verified as far as we know. Guo's
argument is elementary --- Dirichlet's hyperbola method and positivity of
character sums, with rational-inequality numerics --- and is therefore
independently re-checkable, but a reader who declines to rely on it can
substitute Kimura's grid constants into the same computation and recover
$d < e^{19.06} \approx 1.9\times 10^8$; the Corollary additionally relies
on Guo's Appendix~A. Third, our certification is interval arithmetic, not
a Lean kernel check; formalising the arithmetic step would be routine
but has not been done.

\section*{AI-assistance disclosure}

This note was produced by an AI assistant (Ozzie; Muse, Meta),
working under the direction of George, who commissioned the work in
response to Tao's challenge \cite{tao}. The literature search, the
parameter optimisation, the interval-arithmetic certification and the
drafting were all performed by the AI system; all intermediate artefacts
(the computation script, its outputs, and a step-by-step derivation) are
distributed with the note so that every claim can be checked without
trusting the AI. The inputs of Guo \cite{guo} and Kimura \cite{kimura},
and the prior extractions of Rudolph and Murphy reported in \cite{tao},
are used with attribution, as the challenge's norms require.

\section*{Code availability}

\texttt{compute\_bound.py}, its full output
(\texttt{computation\_output.txt} / \texttt{.json}), a term-by-term
breakdown (\texttt{breakdown.txt}) and a prose derivation
(\texttt{DERIVATION.md}) accompany this note in the same folder.

\begin{thebibliography}{9}
\bibitem{tao} T.~Tao, \emph{An experimental AI-math challenge: using
tenth discriminant bounds to explore arguments ruling out Siegel zeroes},
What's New, 20 October 2026. (Challenge post; participant bounds cited
above --- Rudolph, Murphy, twinheterodox, LemmaLab, Larsson, Chung and
the rStation agents --- appear in its comments, 20--28 October 2026.)

\bibitem{openai} OpenAI, \emph{Uniform exclusion of Landau--Siegel
zeros}, 1 October 2026.

\bibitem{banks} W.~D. Banks, \emph{On the OpenAI proof that
Landau--Siegel zeros do not exist}, arXiv:2610.15483.

\bibitem{kerr} B.~Kerr, \emph{A simpler proof that there are no
Landau--Siegel zeros}, arXiv:2610.21099.

\bibitem{green} B.~Green, \emph{On Landau--Siegel zeros},
hexagon:2610.00311.

\bibitem{kimura} I.~Kimura, \emph{Imaginary quadratic fields of class
number at most 2000}, arXiv:2610.21511 (v2).

\bibitem{guo} V.~Z. Guo, \emph{An explicit elementary proof of no
Landau--Siegel zeros}, arXiv:2610.30476 (v1).
\end{thebibliography}

\end{document}
